\section{Encoding the diagrams and their actions}\label{sec:encoding}

The ring needed for the action theorem must encode an action built
from its own elementary group $\E_7(B)$ on $B^7$.  Starting with one
ring and encoding its action in a larger ring would not meet this
requirement: the acting group would still come from the original ring.
We construct one finitely presented ring $B$ that accommodates this
action, the two idempotent diagrams, and a centralizing binary corner.

We first describe the diagrams for an arbitrary acting pair, so the
required ring data can be stated precisely.  The construction then
uses finite codes for ring expressions and action data before their
equality algorithms have been proved to terminate.  A calculation
on a code word of length $n$ will call the prospective normalizer
only on words of strictly smaller length.  Induction on that length
will prove these algorithms total and justify the required ring data.

\subsection{Two enlargements of an acting pair}

Let a group $V$ act on a set $Z$, and let $\ell\in\{2,3\}$.  Use
the prefix algebra $D_\ell$ and finite collection $\mathcal P_\ell$
of Lemma~\ref{lem:idempotents}.  Define
\begin{equation}\label{eq:enlargement-map}
 j_E(g)=(1-E)\otimes1+E\otimes g
 \quad(g\in V,\ E\in\mathcal P_\ell)
\end{equation}
in $D_\ell\otimes_{\F_\ell}\F_\ell[V]$, and let
$\mathcal W_\ell(V)$ be the subgroup of units generated by these
elements.  It acts on
\[
 \mathcal Y_\ell(Z)=D_\ell^{(Z)},
\]
the set of finitely supported arrays with entries in $D_\ell$:
$d\otimes g$ multiplies coefficients on the left by $d$ and sends
index $z$ to $gz$.  Write $\mathbf e_z$ for the unit basis entry.

\begin{lemma}\label{lem:enlargement}
Each $j_E$ is a homomorphism.  The map $j_I$ is injective, and
together with $z\mapsto\mathbf e_z$ it embeds the acting pair.
For every indicated orthogonal sum $E+P$ in $\mathcal P_\ell$,
the images of $j_E$ and $j_P$ commute elementwise and
\[
 j_{E+P}(g)=j_E(g)j_P(g).
\]
If $gz=z$, then $j_E(g)\mathbf e_z=\mathbf e_z$ for every $E$.
\end{lemma}

\begin{proof}
Orthogonality of the two summands in
Equation~\eqref{eq:enlargement-map} gives
$j_E(g)j_E(h)=j_E(gh)$ and inverse $j_E(g^{-1})$.
If $EP=PE=0$, direct multiplication gives the claimed commuting
product.  Since $D_\ell$ is nonzero, the elements $1\otimes g$
for distinct $g$ are distinct group-basis elements.  The action
assertions follow from
\[
 j_E(g)\mathbf e_z=(1-E)\mathbf e_z+E\mathbf e_{gz}.
\]
\end{proof}

For a unital $\F_2$-algebra $B$, put $Q=\E_7(B)$ and $X=B^7$,
with the usual left action on columns.  Apply the enlargement twice:
\begin{equation}\label{eq:two-enlargements}
 (W_1,Y^{(1)})=\bigl(\mathcal W_2(Q),\mathcal Y_2(X)\bigr),
 \qquad
 (W,Y^{(2)})=\bigl(\mathcal W_3(W_1),\mathcal Y_3(Y^{(1)})\bigr).
\end{equation}
Write
\[
 \iota:Q\longrightarrow W,\qquad
 \iota_X:X\longrightarrow Y^{(2)},\qquad
 Y=W\iota_X(X)
\]
for the full-identity embeddings and the indicated invariant subset.
With superscripts indicating the enlargement layer, set
\begin{equation}\label{eq:two-diagrams}
 j_{2,E}=j_I^{(3)}\circ j_E^{(2)},\qquad
 j_{3,E}=j_E^{(3)}\circ j_I^{(2)}.
\end{equation}
Both full-identity maps are $\iota$.  Each family satisfies the
orthogonal identities, and
\begin{equation}\label{eq:diagram-fixed-points}
 qx=x\quad\Longrightarrow\quad
 j_{\ell,E}(q)\iota_X(x)=\iota_X(x).
\end{equation}
The characteristic-three algebra in the second enlargement is used
only to define an abstract group.  We will embed that group, not
its algebra, into the characteristic-two ring.

\begin{theorem}\label{thm:ring-data}
For every finitely generated group $G$ with decidable word problem,
there is a nonzero finitely presented unital $\F_2$-algebra $B$
with the following properties.
\begin{enumerate}[label=\textup{(\roman*)}]
\item There is an injective nonunital algebra endomorphism
  $\phi:B\to B$.  With $p=\phi(1)$ there are
  $s_0,s_1,d_0,d_1\in pBp$, centralizing $\phi(B)$, such that
  \begin{equation}\label{eq:ring-binary-pairs}
   d_i s_j=\delta_{ij}p,\qquad s_0d_0+s_1d_1=p.
  \end{equation}
\item Form $Q,X,W,Y,\iota,\iota_X$ and $j_{\ell,E}$ as above.
  There is a nonzero idempotent $e\in B$ such that $G$ and $W$
  embed in $(eBe)^\times$.  There are $S_y,T_y\in eBe$, for
  all $y\in Y$, with
  \begin{equation}\label{eq:ring-prefix-data}
   T_yS_z=\delta_{yz}e,\qquad
   wS_y=S_{wy},\qquad T_yw=T_{w^{-1}y}.
  \end{equation}
  Here $w$ denotes its corner-unit image.
\end{enumerate}
The two diagrams have common full-identity map $\iota$ and
satisfy Equation~\eqref{eq:diagram-fixed-points}.
\end{theorem}

The construction has four steps. We first recall the compiler and
binary-corner results that we use. We then describe a monoid encoding
a group action and give finite codes for the prospective ring and
action data. A length induction proves that the resulting
self-referential normalizer is total and consistent in contexts.
Finally, contracted monoid algebras turn the codes into the required
ring embeddings and corner relations.

\subsection{The bounded compiler and the auxiliary monoid}

We use the following established form of the finite compiler
\citep[Lemma~7.1, Bounded Compiler Lemma]{FoundationsCompanion}.
Its bounded contextual
conclusion is essential.

\begin{lemma}[Bounded compiler]\label{lem:compiler}
Fix a sufficiently large finite alphabet $\Sigma$ containing the
distinguished letters of the compiler.  From an index $\eta$ for
a partial function $f_\eta:\Sigma^*\rightharpoonup\Sigma^*$ one
can compute a finite monoid presentation $K_\eta$, with alphabet
$\mathcal A_\eta$ and zero letter $\Omega$, and an additive
nonnegative integral word grade
$\gamma:\mathcal A_\eta^*\to\Z_{\geq0}$. Letter grades are at
most two, and $\gamma(\Omega)=0$.

There is a partial algorithm
$N_\eta:\mathcal A_\eta^*\rightharpoonup\mathcal A_\eta^*$,
independent of any budget,
with the following property.  Suppose that, for some $m\geq0$,
\begin{align}
 &f_\eta(u)\text{ is defined and }|f_\eta(u)|\leq |u|
       &&(|u|\leq m),\label{eq:compiler-budget-total}\\
 &f_\eta(a f_\eta(u)b)=f_\eta(aub)
       &&(|aub|\leq m).\label{eq:compiler-budget-context}
\end{align}
Then $N_\eta$ terminates through grade $m$.  Each such input is
connected to its literal output by a defining-equation derivation
that never raises grade.  A contextual defining-equation step
preserves the output if both endpoints have grade at most $m$.
Consequently,
\begin{equation}\label{eq:compiler-context}
 N_\eta(u)=N_\eta(v)\ \Longrightarrow\
 N_\eta(aub)=N_\eta(avb)
 \quad\text{if }\gamma(aub),\gamma(avb)\leq m.
\end{equation}
Unconditionally, $N_\eta$ fixes the empty word and $\Omega$,
sends every word containing $\Omega$ to $\Omega$, and fixes a
specified literal word $\operatorname{raw}(u)$ for every
$u\in\Sigma^*$.  The raw words are mutually distinct and distinct
from the empty word and $\Omega$.  If the hypotheses hold at every
$m$, then $N_\eta$ decides equality in $K_\eta$ and
\begin{equation}\label{eq:compiler-raw-product}
 N_\eta\bigl(\operatorname{raw}(u)\operatorname{raw}(v)\bigr)
   =\operatorname{raw}\bigl(f_\eta(uv)\bigr).
\end{equation}
The presentation is computed without evaluating $f_\eta$.
Apart from finite computations, $N_\eta$ calls $f_\eta$ only on
inputs of length at most the grade of its own input.
\end{lemma}

Fix also the monoid $\mathcal D$ of partial maps of $\Q$ generated
by the full affine maps
\[
 u(x)=x/2,\qquad v(x)=(x+1)/2,
\]
their inverses, and the partial identities on
$[0,\infty)\cap\Q$ and $[1,\infty)\cap\Q$.
Composition applies the right factor first.  An element has
computable data $(a,b,r)$, denoting $x\mapsto ax+b$ on
$[r,\infty)\cap\Q$, with $a=2^j$ for $j\in\Z$, dyadic $b,r$,
and $r=-\infty$ for full domain.  These data multiply by
\begin{equation}\label{eq:partial-affine-data}
 (a,b,r)(c,d,s)
   =\bigl(ac,ad+b,\max\{s,(r-d)/c\}\bigr).
\end{equation}
Every domain is infinite, so equality is equality of these data.
This gives an unconditional decision algorithm.

We will use the precise conclusion of the Binary Tail Corner Lemma
\citep[Lemma~6.4]{FoundationsCompanion} for this same monoid:
\begin{equation}\label{eq:external-binary-corner}
 \begin{gathered}
 0\ne q=q^2\in\F_2[\mathcal D],\qquad
 s'_i,d'_i\in q\F_2[\mathcal D]q,\\
 d'_is'_j=\delta_{ij}q,\qquad
 s'_0d'_0+s'_1d'_1=q.
 \end{gathered}
\end{equation}
Its $q$ is the difference of the two tail identities, split
using the affine maps $u$ and $v$.

\subsection{A monoid for an action}

For an acting pair $V\curvearrowright Z$, define $M(V,Z)$ with
absorbing zero, generators $V$ and $\mathsf s_z,\mathsf t_z$
for $z\in Z$, and relations consisting of group multiplication,
zero absorption, and
\begin{equation}\label{eq:action-monoid-relations}
 \mathsf t_y\mathsf s_z=\delta_{yz},\qquad
 g\mathsf s_z=\mathsf s_{gz},\qquad
 \mathsf t_zg=\mathsf t_{g^{-1}z}.
\end{equation}
Here $\delta_{yz}$ denotes the monoid identity or zero.

\begin{lemma}\label{lem:action-monoid-normal-form}
Every nonzero element of $M(V,Z)$ has a unique normal form
\[
 \mathsf s_{z_1}\cdots\mathsf s_{z_a}\,g\,
 \mathsf t_{w_1}\cdots\mathsf t_{w_b},
 \qquad a,b\geq0,
\]
omitting an identity group label.  In particular, $V$ embeds
as a group of units.
\end{lemma}

\begin{proof}
Multiply adjacent group labels, erase identity labels, and
orient Equation~\eqref{eq:action-monoid-relations} to the right.
Each rule shortens the word.  Overlapping group products agree
by associativity; overlaps of two group labels with a prefix
or deletion agree by the action law.  The remaining nonzero
overlap is $\mathsf t_y g\mathsf s_z$.  Its two reductions
agree because $y=gz$ if and only if $g^{-1}y=z$.
Identity and zero overlaps join immediately.  The system is
therefore terminating and locally confluent, hence confluent.
Its irreducibles are exactly the stated words.
\end{proof}

\subsection{Syntax and costs before totality}

For a field $k$ and a monoid $M$ with absorbing zero, write $k_0[M]$ for its
\emph{contracted monoid algebra}: its basis consists of the nonzero
elements of $M$, and a monoid product equal to zero is interpreted
as algebra zero.

Fix provisionally an arbitrary index $\eta$ and its compiler
presentation.  Polynomial expressions over $\F_2$ are finite
sums of monoid words in the presentation letters.  Their
\emph{cost} is the maximum grade of a monomial, with cost zero
for the zero expression.  A column expression has the maximum
cost of its seven entries.  These are syntactic expressions:
we do not assume global equality in the prospective contracted
algebra $(\F_2)_0[K_\eta]$.

Use three families of generating symbols.
\begin{enumerate}[label=\textup{(\alph*)}]
\item A finite symmetric generating alphabet for $G$, evaluated
  using its word decider.  These symbols have cost zero.
\item The presentation alphabet $\mathcal A_\eta$ and a fixed
  finite generating alphabet for $\mathcal D$.  Their prospective
  values lie in the zero-product monoid
  \begin{equation}\label{eq:zero-product-monoid}
   T_\eta=\bigl((K_\eta\setminus\{\Omega\})\times\mathcal D\bigr)
          \sqcup\{0\}.
  \end{equation}
  Nonzero pairs multiply coordinatewise, with a product having
  first coordinate $\Omega$ replaced by zero.  Presentation
  letters have their compiler grades as costs; the other
  symbols have cost zero.
\item Root symbols $I+b\mathbf e_{ij}$, for $1\leq i\ne j\leq7$,
  where $\mathbf e_{ij}$ is the matrix unit and $b$ is an arbitrary
  polynomial expression, and their formal inverses.
  Carry each root through $j_F^{(3)}j_E^{(2)}$, with
  $E\in\mathcal P_2$ and $F\in\mathcal P_3$.
  Also use $\mathsf s_x,\mathsf t_x$ for every column expression
  $x$, indexed by its twice embedded basis entry, and a zero
  symbol.  Costs are those of the root entry or column;
  zero has cost zero.  The prospective evaluation is in $M(W,Y)$.
\end{enumerate}
These alphabets are effectively enumerable from $\eta$ without
running $N_\eta$.  Distinct syntactic descriptions may have the
same eventual value.  The third alphabet generates $M(W,Y)$
once the prospective objects exist: roots generate $Q$, each
$j_E$ is a homomorphism, and translating the supplied column
symbols gives those indexed by all of $Y=W\iota_X(X)$.

Encode these symbols by distinct computably decodable words in
a fixed finite alphabet $\Sigma$, reserving a letter $\zeta$ for
zero. The families use disjoint alphabets, and none uses $\zeta$.
Within each family, fix a start marker and a different end marker.
Every code in that family begins and ends with these markers,
neither of which occurs internally.  Binary syntax descriptions with unary
padding allow us to require
\begin{equation}\label{eq:code-cost}
 |\operatorname{code}(a)|\geq10\operatorname{cost}(a)+2.
\end{equation}
Fix a total order on $\Sigma$. The \emph{shortlex} order compares
words first by length and then lexicographically.
Code membership and enumeration of all code concatenations
through any specified length require no call to $N_\eta$.
The cost of a sequence of symbols is the sum of their costs.

\subsection{Calculations in a bounded system}

The code normalizer will compare expressions and replace them by
shorter codes for the same value.  A shorter code can have greater
symbolic cost.  We therefore need equality to survive substitution
whenever each resulting expression separately fits the available
budget.  The following lemma establishes that property without
assuming that the prospective ring or groups have global equality
algorithms.

\begin{lemma}\label{lem:bounded-action-calculus}
Suppose the hypotheses of Lemma~\ref{lem:compiler} hold through
length $3d$, where $d\geq0$. Within any one of the three families,
evaluation of any symbol sequence of
cost at most $d$, and equality comparison of two such sequences
(each separately of cost at most $d$), terminate using compiler
calls of grade at most $3d$. Tested equality respects
replacing a subword in a product whenever both whole symbol
sequences have cost at most $d$.  Tests agree at different
valid budgets.
\end{lemma}

\begin{proof}
Put $m=3d$.  We first construct bounded algebra, group, and action
operations.  We then normalize words in the action monoid and prove
that equal normal forms can be substituted even when their original
representatives have different costs.

\paragraph{Bounded algebra and action operations.}
For $0\leq b\leq m$, form the classes represented
by words of grade at most $b$, where equality is literal
equality of $N_\eta$ outputs.  These sets are nested, since
equality does not depend on the budget.  Products are defined
when the summed grades fit $m$.  By
Equation~\eqref{eq:compiler-context}, changing representatives
preserves a product whenever both resulting grades fit.
The nonincreasing normalization derivation puts the output
in the ball of the summed grade.  Associativity follows by
concatenating three representatives within their total budget.
The empty and zero classes are distinct, and zero absorbs.

A class has a minimum represented grade by well-ordering.
We use this minimum only in proofs, never as an algorithmic
search.  Finite $\F_2$-linear combinations of nonzero classes
give nested algebra balls.  Addition takes the maximum input
bound, multiplication the sum.  Equality is comparison of
basis coefficients after collecting equal $N_\eta$ outputs.
Distribution and collection preserve products by bounded
contextual compatibility.

For a word in elementary root matrices of cost at most $b$,
store both its forward matrix and the matrix of its formal
inverse word.  Their entries belong to the algebra ball of
level $b$.  Compare both components.  Operations are
\begin{equation}\label{eq:bounded-matrix-pairs}
 (A,A^-)(C,C^-)=(AC,C^-A^-),\qquad
 (A,A^-)^{-1}=(A^-,A).
\end{equation}
They respect equality whenever the summed levels fit $m$.
This avoids assuming uniqueness of inverses outside the
available budget.  A word and its formal inverse cancel
when their total cost fits, by successively cancelling
adjacent root factors; intermediate calculations stay in
that total bound.  Forward matrices act on columns, with
group bound $b$ and column bound $c$ giving bound $b+c$.

This construction propagates through an enlargement.
Expand a group word into a finite array of $D_\ell$-coefficients
indexed by preceding-level group labels; separately expand
the formal inverse.  For a word of cost at most $b$, all
indices in both arrays have level at most $b$.  Compare
arrays after collecting equal indices and performing the
unconditional arithmetic in $D_\ell$.  Multiplication
convolves arrays and adds bounds, reversing the order in
the inverse component as in
Equation~\eqref{eq:bounded-matrix-pairs}.

The corresponding set ball of level $c$ consists of finite
arrays with preceding-level set indices of level at most
$c$.  An acting array of level $b$ gives an output of level
$b+c$, by multiplying coefficients and acting on indices.
Expanding a full triple product gives identical arrays in
either association if the total input bound is at most $m$.
The same expansion proves the action law.  The maps $j_E$
and the basis-entry maps propagate generators and points
without increasing cost.  Formal inverse cancellation
follows from $E^2=E$ and the preceding inverse identities.
Repeating this argument gives both enlargement levels.
Their level inclusions are injective because equality is
literal equality of recursively collected arrays.

We have now defined products, inverses, and actions at both
enlargement levels whenever the total input costs fit $m$.
All identities used to define these operations have been checked
within that budget.  The algorithm retains the original words and
polynomials, multiplies by concatenation and finite expansion, and
uses $N_\eta$ for comparison and collection.  Thus these are finite
calculations in truncated systems; no complete group is assumed,
and no minimum-cost representative is sought.

\paragraph{Normal forms for the action monoid.}
Apply the reductions of
Lemma~\ref{lem:action-monoid-normal-form} to the third
family.  Every reduction consumes disjoint input data,
so the output labels have representations of jointly
bounded total cost no greater than the input cost.
Associativity and the action law resolve their usual
overlaps.  For the exceptional overlap
$\mathsf t_y w\mathsf s_z$, let the costs of $y,w,z$
be $c,b,e$, with $b+c+e\leq d$.  The equivalence
\begin{equation}\label{eq:bounded-inverse-action}
 y=wz\quad\Longleftrightarrow\quad w^{-1}y=z
\end{equation}
uses at most $2b+e$ or $2b+c$ when acting back by the
inverse or forward element.  Both are at most $2d$,
hence within $m$.  For example, equality $y=wz$
allows the action on either representative, and
$w^{-1}(wz)=z$ by the bounded inverse identity.
The converse is the same calculation.  Identity
and zero overlaps join immediately.  Thus the
terminating reductions are confluent on inputs
of total cost at most $d$, and their normal forms
are computed by the finite operations just described.

\paragraph{Substitution with unequal costs.}
It remains to prove the contextual assertion.  The difficulty is
that tested-equal subwords may have different costs, so one cannot
charge both representatives to a single calculation.
Give each group or point label its minimum represented cost in the
truncated system.  At every group level, take this minimum only over
genuine words in that level's generators, stored together with their
formal inverses.  Point labels may instead use ambient array
representatives.

Consider abstract labeled words whose total minimum cost is at most
$d$.  Each reduction above does not increase that total: the resulting
labels have representatives whose joint cost is no larger than that
of the labels consumed.  The same overlap calculations therefore
prove confluence on these abstract words.  In particular, the
inverse-action overlap still fits $2d\leq m$.

By the equality test just constructed, two tested-equal subwords
have the same normal form.  If it is zero, both contextual products
reduce to zero.  Otherwise it is a list $R$ of prefix labels, group
label, and deletion labels.  Its total minimum
cost is at most the original cost of either subword, since each
reduction supplied jointly bounded representatives for its output.
Put the subwords in fixed left and right contexts, with each whole
word separately of cost at most $d$.  Normalize the middle first.
Both routes reach the identical abstract word with middle $R$, and
this word fits the budget on each route.  Its unique normal form is
therefore the answer on both sides.  No step adds the costs of the
old and new equal representatives.

Minimum-cost point representatives may be ambient
array expressions rather than short orbit words.
That is sufficient: $R$ is a semantic proof device,
and the algorithm never needs to find a code for
it or decide membership in an orbit.

The second family has contextual compatibility
directly from Equation~\eqref{eq:compiler-context}
and the affine data; the first has it from its
group law and word decider.  Finally all the
comparisons use outputs of the same algorithm
$N_\eta$, not a budget-dependent normalizer.
The recursive array comparisons inherit that
independence, proving compatibility of budgets.
\end{proof}

\subsection{The self-referential normalizer}

The bounded calculations are now available conditionally.
We next construct the word normalizer whose consistency is the
condition they require. Its call bounds will make this apparent
circularity into an induction on word length.

For $\eta$ and a word $w\in\Sigma^*$ of length $n$,
define a partial procedure $F(\eta,w)$ as follows.
Include $a\zeta\to\zeta$ and $\zeta a\to\zeta$ for
every $a\in\Sigma$.  For every nonempty concatenation
$c$ of whole codes from one family with $|c|\leq n$,
evaluate its prospective product.  If the value is
zero, include $c\to\zeta$.  Otherwise search all
same-family code concatenations of length at most
$|c|$, also allowing the empty word for identity,
and choose the least tested-equal one in shortlex.
Include the replacement unless unchanged.  The
starting word is itself a candidate, so the search
is finite whenever its tests terminate.

Reduce $w$ using a fixed deterministic choice of
rule and occurrence.  Every nontrivial rule decreases
shortlex.  A zero-valued code has length at least
two, so replacement by $\zeta$ also shortens.
Thus the reduction terminates once its finite
table of tests is available.  This is uniformly
partial computable, including computation of the
compiler presentation from its parameter $\eta$.

Effective parameterization gives a total computable
index transformation $\tau(\eta)$ for $F(\eta,-)$.
Kleene's Recursion Theorem supplies an index such that
\begin{equation}\label{eq:encoding-fixed-point}
 f_\eta(w)=F(\eta,w)
\end{equation}
as partial functions \citep[\S~2, p.~153]{Kleene38}.
Explicitly, if $U$ is universal evaluation, let
$d(a)$ be an effectively obtained index for
$w\mapsto U(U(a,a),w)$.  Choose $b$ computing
the total index-valued function $a\mapsto\tau(d(a))$.
Then $\eta=d(b)$ satisfies
\[
 U(\eta,w)=U(U(b,b),w)
          =U(\tau(d(b)),w)=F(\eta,w).
\]
No termination of $F(\eta,-)$ is assumed here.
Fix this $\eta$.

\begin{lemma}\label{lem:encoding-bootstrap}
The function $f_\eta$ is total and length-nonincreasing.
Its rule tables are stable on bounded-length words and
confluent on those words.  In particular,
\begin{equation}\label{eq:encoding-context-law}
 f_\eta(a f_\eta(u)b)=f_\eta(aub)
 \qquad(a,u,b\in\Sigma^*).
\end{equation}
\end{lemma}

\begin{proof}
Induct simultaneously on the input length $n$.
For $n=0$ there are no tests or reductions and
the empty word is unchanged.  For $n>0$, put
$d=\lfloor n/10\rfloor$.  Every code string in
a table entry or representative search has cost
at most $d$ by Equation~\eqref{eq:code-cost}.
Since
\begin{equation}\label{eq:strict-encoding-dependency}
 3\lfloor n/10\rfloor<n,
\end{equation}
the induction hypothesis supplies the compiler
hypotheses through length $3d$.  All tests
terminate by Lemma~\ref{lem:bounded-action-calculus}.
Only finitely many code strings have length at
most $n$, even though the symbol alphabet is
countable and can contain infinitely many
zero-cost symbols.  The finite table and its
shortlex-decreasing reduction therefore terminate.

A rule on a fixed code string is determined by
the same tests and the same search bound at
every larger length.  Thus rules with left
side of length at most $k$ are unchanged once
length $k$ is reached; longer rules cannot
act on a word of length at most $k$.

For local confluence in the length-$n$ ball,
disjoint occurrences commute, even when a
replacement is empty.  Every branch from a
word containing $\zeta$ retains it and can
be absorbed to $\zeta$, since code strings
avoid that letter.  Two overlapping code
rules have aligned code boundaries and
belong to the same family, by the delimiter
conditions.  Their union is a whole code
concatenation of length at most $n$.

Either replacement preserves its tested class
by Lemma~\ref{lem:bounded-action-calculus}.
A shorter code representative may have higher
symbolic cost.  We do not assert otherwise:
the old and new unions each separately have
length at most $n$, hence cost at most $d$.
This is the two-endpoint hypothesis of that lemma.

For a zero-class union, a branch either
introduces $\zeta$ and absorbs, or leaves
a zero-class code string with a whole-word
zero rule.  It cannot leave the empty
string, since zero and identity are distinct.
For a nonzero-class union, let $c_*$ be
its least representative within the original
search bound.  A branch gives a same-class
string $c_b$ of no greater length.  This was
an original candidate, so $c_*\leq c_b$
in shortlex, in particular $|c_*|\leq|c_b|$.
Thus $c_*$ is still the least candidate in
the branch's smaller search bound.  Each
branch reaches it by its whole-word rule,
unless already there.  This covers an empty
$c_*$ and containment overlaps as well.

The joining steps occur inside the old
union before any newly adjacent outside
letters are used.  Erasing identity codes
can therefore expose further redexes but
creates no missing local peak.  In particular,
different-family adjacencies after an erasure
do not create an overlapping initial peak.

Termination and local confluence give
confluence by Newman's Lemma \citep{Newman42}.
Its well-founded proof applies within the
ball: proper descendants have unique normal
forms, and a common descendant of two first
steps equates their normal forms.

For $|aub|\leq n$, normalize $u$ inside
its context first.  Stability puts all
these rules in the length-$n$ table;
length never increases.  Confluence proves
$f_\eta(a f_\eta(u)b)=f_\eta(aub)$.
The compiler hypotheses now hold through
length $n$, completing the induction.
\end{proof}

\subsection{The monoids and the ring}

The length induction has supplied a total normalizer consistent in
every context.  We now obtain the two algebra embeddings used to
construct $B$ and its self-map.  One comes directly from the code
families.  For the other, the compiler's raw words preserve monoid
multiplication but send the absorbing zero to a nonzero algebra
basis element.  Subtracting that basis element will give an embedding
of the contracted monoid algebra.

Let $L$ be the monoid on $\Sigma$ with equations
the union of the rule tables.  Every word is
related to its computed normal form.  Conversely,
a contextual defining equation preserves that
form: choose a length bound containing the
equation and both endpoints, and use stability
and confluence.  Thus equality in $L$ is exactly
equality of literal $f_\eta$ outputs.  Its
identity and absorbing zero are the distinct
irreducibles $\varepsilon$ and $\zeta$.

Lemma~\ref{lem:compiler} now applies at every
grade, giving an equality algorithm for
$K=K_\eta$ with distinct identity and zero.
Set
\[
 k=\F_2,\qquad B=k_0[K],\qquad A_*=k_0[L].
\]
The previously prospective $Q,X,W,Y$
now exist, and the bounded calculations
are their genuine evaluations.  Paired
forward/inverse equality agrees with
ordinary group equality, since the
formal inverse is now a genuine inverse.

\begin{lemma}\label{lem:encoding-embeddings}
The code families embed $G$, $T_\eta$, and
$M(W,Y)$ unitally into $L$, preserving
zeros in the last two cases.  There are
injective algebra maps
\begin{equation}\label{eq:encoding-algebra-maps}
 \kappa:B\otimes_k k[\mathcal D]\longrightarrow A_*,
 \qquad i:A_*\longrightarrow B,
\end{equation}
where $\kappa$ is unital and $i$ has a
nonzero proper support identity $e=i(1)$.
\end{lemma}

\begin{proof}
In one family, a least shortlex code
representative of a nonzero value is
irreducible.  An internal code reduction
would preserve its value and lower
shortlex, while zero and other-family
rules cannot occur inside it.  Every other
representative reduces to it by its
whole-word rule.  Zero values reduce
to $\zeta$.  Hence values remain distinct.
Codes concatenate and the empty identity
representative is allowed, giving unital
multiplicative embeddings.  The embedded
groups are units, using also
Lemma~\ref{lem:action-monoid-normal-form}.

For $h\in L$ with normal word $\overline h$, put
\[
 c(h)=[\operatorname{raw}(\overline h)]\in K.
\]
Distinct normal words give distinct
literal raw $N_\eta$ outputs, hence
distinct elements of $K$ avoiding
$1_K$ and $\Omega$.
Equation~\eqref{eq:compiler-raw-product}
gives $c(h)c(h')=c(hh')$.
The identity and zero need not be
preserved: $c(1_L)\ne1_K$ and
$c(\zeta)\ne\Omega$.  They are identity
and absorbing zero only within $c(L)$.

For a nonzero basis element $h$ of $A_*$, set
\begin{equation}\label{eq:raw-difference-embedding}
 i(h)=c(h)-c(\zeta).
\end{equation}
Since $c(h)c(\zeta)=c(\zeta)c(h)=c(\zeta)$
and $c(\zeta)^2=c(\zeta)$,
\[
 (c(h)-c(\zeta))(c(h')-c(\zeta))
     =c(hh')-c(\zeta).
\]
If $hh'=\zeta$, this is zero as required.
Distinct $c(h)$ are distinct nonzero
basis vectors of $B$.  The vectors in
Equation~\eqref{eq:raw-difference-embedding}
are therefore linearly independent,
proving injectivity.  The support identity
\[
 e=c(1_L)-c(\zeta),\qquad e^2=e,
\]
is nonzero and is not $1_B$, since neither
of its basis vectors is $1_K$.

Finally nonzero basis pairs identify
$B\otimes_k k[\mathcal D]$ with
$k_0[T_\eta]$.  The code embedding
of $T_\eta$ extends linearly to
the required unital injection $\kappa$.
\end{proof}

\begin{lemma}\label{lem:encoding-finite-presentation}
The algebra $B$ is nonzero and finitely
presented as a unital associative
$\F_2$-algebra, and as a unital ring.
\end{lemma}

\begin{proof}
Take the finite generators of $K$ in
the free unital associative $\F_2$-algebra.
Replace its defining equations $u=v$
by binomials $u-v=0$, and impose
$\Omega=0$.  Before the last relation,
the binomial ideal is exactly the span
of differences of words in the same
monoid congruence class.  Indeed
contextual equation steps lie in the
ideal and telescope along a derivation;
conversely the linear map sending words
to monoid classes annihilates all
binomials.  The quotient has basis $K$.
Killing its absorbing zero class removes
exactly that basis element and gives
$k_0[K]$.  Identity survives since
$1_K\ne\Omega$.  Adding $2\cdot1=0$
gives a presentation as a unital ring.
\end{proof}

\begin{proof}[Proof of Theorem~\ref{thm:ring-data}]
The ring properties follow from
Lemma~\ref{lem:encoding-finite-presentation}.
The code embeddings put $G$ and $W$
in $A_*^\times$.  Transporting them
through $i$ puts them in $(eBe)^\times$.
Transport the prefix and deletion
symbols to $S_y,T_y\in eBe$.
Their monoid relations become exactly
Equation~\eqref{eq:ring-prefix-data},
with identity mapped to $e$ and zero
mapped to algebra zero.

Use Equation~\eqref{eq:external-binary-corner}
and define
\begin{equation}\label{eq:encoding-phi}
 \phi(b)=i\kappa(b\otimes q),\qquad
 s_i=i\kappa(1\otimes s'_i),\qquad
 d_i=i\kappa(1\otimes d'_i).
\end{equation}
The first map is multiplicative because
$q^2=q$ and injective because $q\ne0$,
the tensor product is over a field,
and $i,\kappa$ are injective.  With
$p=\phi(1)$, the binary identities
give $s_i,d_i\in pBp$ and
Equation~\eqref{eq:ring-binary-pairs}.
Moreover
\[
 (b\otimes q)(1\otimes s'_i)
   =b\otimes s'_i
   =(1\otimes s'_i)(b\otimes q),
\]
and the same holds for $d'_i$.
Thus these elements centralize $\phi(B)$.
Since $p\in eBe$ and $e\ne1_B$,
we have $p\ne1_B$, so $\phi$ is
nonunital.  The diagram assertions
follow from Lemma~\ref{lem:enlargement}
and Equations~\eqref{eq:two-diagrams}
and~\eqref{eq:diagram-fixed-points}.
\end{proof}
